\section{进入汽车大本营：从样机到主机厂体系}

上一章讲客户倒戈，本章讲进入汽车大本营。汽车产业链对供应商要求极高：车规、可靠性、认证、成本、量产、售后、长期供货都必须过关。激光雷达从样机到主机厂体系，不是“卖一个硬件”，而是成为整车安全和智能系统的一部分。

\begin{figure}[H]
\centering
\includegraphics[width=0.96\textwidth]{auto-maincamp.png}
\caption{进入汽车大本营：从技术样机到车规供应商，要进入主机厂体系。自制概念图，依据 01:58:07--02:38:34 对谈内容整理。}
\end{figure}

\begin{knowledgebox}{读图：主机厂体系为什么难}
主机厂要的是长期稳定供应，而不是一次性技术惊艳。供应商必须证明产品可靠、成本可控、交付稳定、售后可承担，并能随着车型周期持续迭代。
\end{knowledgebox}

\subsection{汽车产业链机会来自哪里}

上一节讲主机厂体系为什么难，本节追问：这样难的产业机会到底来自哪里。李一帆反复谈机会。汽车产业链机会不是单个创业者凭空创造，而是新能源车崛起、供应链国产化、智能化需求、主机厂竞争和国家产业能力共同作用。禾赛恰好站在这个多重机会交汇处。

\begin{importantbox}{产业机会的结构}
\[
\text{Opportunity} \approx \text{Technology} \times \text{Supply Chain} \times \text{Market Timing} \times \text{Customer Adoption}
\]
其中，Technology 是技术成熟，Supply Chain 是供应链可实现，Market Timing 是产业窗口，Customer Adoption 是主机厂愿意采用。
\end{importantbox}

\begin{knowledgebox}{汽车产业链机会分解}
\begin{tabularx}{\textwidth}{p{0.20\textwidth}p{0.34\textwidth}X}
\toprule
机会来源 & 表现 & 对激光雷达公司的意义 \\
\midrule
新能源车增长 & 新车型、新平台、新供应链 & 有机会进入下一代车型设计。 \\
智能化竞争 & 辅助驾驶、Robotaxi、机器人需求上升 & 感知传感器价值被重新评估。 \\
国产供应链 & 主机厂愿意寻找本土供应商 & 给新公司切入窗口。 \\
成本下降 & 高端配置下探到更多车型 & 市场规模扩大。 \\
全球化需求 & 海外客户需要替代供应链 & 要补国际信任和车规经验。 \\
\bottomrule
\end{tabularx}
\end{knowledgebox}

\begin{importantbox}{从样机公司到车规供应商}
样机公司卖的是可能性，车规供应商卖的是长期确定性。前者需要技术惊艳，后者需要质量体系、供应链、售后、成本曲线和客户信任。
\end{importantbox}

\subsection{本章小结}

进入汽车大本营意味着公司从技术供应商变成产业链节点。真正的竞争从此不只是性能，而是车规、成本、交付和主机厂信任。

